\chapter{Circuit QED: coupling, control, and gates}
\label{ch:circuit-qed}

Parts~I--III built a complete toolkit --- two-level systems, harmonic
oscillators, the Jaynes--Cummings coupling, the dispersive shift, drives
and the RWA --- but left one number, the coupling $g$, as a free
parameter handed down from somewhere. This chapter is that somewhere.
With a concrete qubit (the transmon of \cref{ch:transmon}) and a
concrete resonator (the quantised $LC$ circuit of
\cref{ch:circuit-quantisation}), we can compute $g$ from circuit
capacitances, recover the Jaynes--Cummings model of
\cref{ch:jaynes-cummings} and the dispersive readout of
\cref{ch:dispersive} as hardware facts rather than postulates, and
turn the abstract single- and two-qubit gates of \cref{ch:composite}
into microwave pulses and tunings. The chapter then exploits the one
feature that distinguishes a superconducting circuit from a natural
atom --- its Josephson nonlinearity is \emph{engineerable and
pumpable} --- to generate the parametric and mixing processes behind
quantum-limited amplification, the last link in the readout chain that
\cref{ch:dispersive} left as ``ideally a quantum-limited amplifier''.
The experimental procedures that calibrate all of this follow in
\cref{ch:experimental-basics}.

% --------------------------------------------------------------
\section{Coupling a transmon to a resonator}
\label{sec:cqed-coupling}

\begin{phenomenon}
Place a transmon next to a superconducting resonator and join them with
a small coupling capacitor. The resonator's zero-point voltage
fluctuations push charge onto the transmon island, and the transmon's
charge oscillations drive the resonator: the two share energy. Because
both are, to lowest order, harmonic oscillators
(\cref{stmt:lc-oscillator,stmt:transmon}), their capacitive coupling is
a bilinear exchange interaction --- exactly the atom--field coupling of
cavity QED, now built from lithography rather than a vacuum chamber and
a beam of atoms.
\end{phenomenon}

\begin{statement}[\textbf{Capacitive coupling and the Jaynes--Cummings model}~\cite{blais2004,koch2007}]
\label{stmt:cqed-coupling}
A transmon coupled to a resonator mode through a capacitor $C_g$ has
the interaction
\begin{equation}
\label{eq:cqed-interaction}
\op H_{\rm int} = 2\beta e\,V_{\rm zpf}\,\op n\,(\hat a+\hat a^\dagger),
\qquad
\beta=\frac{C_g}{C_\Sigma},
\qquad
V_{\rm zpf}=\sqrt{\frac{\hbar\omega_r}{2C_r}},
\end{equation}
where $\op n$ is the transmon charge operator, $V_{\rm zpf}$ the
resonator's zero-point voltage, and $\beta$ the fraction of the
resonator voltage that falls across the junction. Writing the transmon
in its eigenbasis, $\op n\approx in_{\rm zpf}(\hat b^\dagger-\hat b)$
with $n_{\rm zpf}\approx(E_J/32E_C)^{1/4}$, and applying the RWA
(\cref{stmt:rwa}) gives the Jaynes--Cummings Hamiltonian
\cref{eq:jc-hamiltonian} of \cref{ch:jaynes-cummings},
\begin{equation}
\label{eq:cqed-jc}
\op H \approx \hbar\omega_r\hat a^\dagger\hat a
+ \frac{\hbar\omega_q}{2}\hat\sigma_z
+ \hbar g\,(\hat a^\dagger\hat\sigma_- + \hat a\hat\sigma_+),
\qquad
\boxed{\,g \propto \frac{C_g}{\sqrt{C_rC_\Sigma}}\sqrt{\omega_r\omega_q}\,}.
\end{equation}
The coupling is fully determined by the circuit: it grows with the
coupling capacitor $C_g$ and with the geometric mean of the two
frequencies, and is set at design time.
\end{statement}

\paragraph{Derivation sketch.}
The two nodes (resonator and qubit island) have charges $Q_r,Q_q$ and
a mutual capacitance $C_g$. The electrostatic energy
$\tfrac12\mathbf Q^\top \mathbf C^{-1}\mathbf Q$ with the capacitance
matrix $\mathbf C$ has an off-diagonal entry $\propto C_g$ that couples
$Q_r$ to $Q_q$. Expressing $Q_r$ through the resonator ladder operators
($Q_r\propto V_{\rm zpf}(\hat a+\hat a^\dagger)$) and $Q_q=2e\op n$
through the transmon ladder operators gives
\cref{eq:cqed-interaction}; keeping only excitation-conserving terms in
the RWA yields \cref{eq:cqed-jc}. Everything beyond this point ---
vacuum Rabi splitting $2g$ (\cref{stmt:vacuum-rabi}), the cQED regimes
and cooperativity (\cref{stmt:cqed-regimes}), and the dispersive shift
$\chi=g^2/\Delta$ (\cref{stmt:dispersive-H}) --- now follows with $g$ a
known number.

\paragraph{The transmon is a multilevel atom.}
One feature has no atomic analogue made this simple: the transmon's
weak anharmonicity $\alpha=-E_C$ (\cref{stmt:transmon}) means its third
level $\ket2$ is only $|\alpha|$ away from harmonic, so it participates
in the dynamics. This is a nuisance for single-qubit gates (leakage)
but a \emph{resource} for two-qubit gates and for parametric processes,
both of which exploit the higher levels directly. We keep the
two-level $\hat\sigma$ notation where it suffices and restore $\ket2$
where it matters.

\paragraph{Transition.}
With $g$ in hand, control is the next ingredient: we drive the qubit
to rotate it, and the resonator to read it. Single-qubit rotations come
first.

% --------------------------------------------------------------
\section{Driving the qubit: single-qubit gates}
\label{sec:single-qubit-gates}

\begin{statement}[\textbf{Drive Hamiltonian and single-qubit rotations}]
\label{stmt:single-qubit-gates}
A microwave tone of envelope $\Omega(t)$, frequency $\omega_d$, and
phase $\phi$, applied to the qubit through a drive capacitor, adds
\begin{equation}
\label{eq:drive-hamiltonian}
\op H_d(t) = \hbar\Omega(t)\cos(\omega_d t+\phi)\,(\hat b+\hat b^\dagger).
\end{equation}
In the frame rotating at $\omega_d$ and on resonance
($\omega_d=\omega_q$), the RWA (\cref{stmt:rwa}) reduces this to
\begin{equation}
\label{eq:rotating-frame-drive}
\op H_d^{\rm rot} = \frac{\hbar\Omega(t)}{2}\bigl(\cos\phi\,\hat\sigma_x + \sin\phi\,\hat\sigma_y\bigr),
\end{equation}
a rotation about an axis in the equatorial plane set by the drive phase
$\phi$, through an angle $\theta=\int\Omega(t)\,\dd t$ set by the pulse
area. Amplitude controls the angle (a $\pi$-pulse inverts the qubit, a
$\pi/2$-pulse creates an equal superposition); phase controls the axis
($\phi=0$ gives $\op R_x$, $\phi=\pi/2$ gives $\op R_y$); and $\op R_z$
is implemented ``virtually'', by advancing the phase of all subsequent
drives. These generate any single-qubit unitary.
\end{statement}

\paragraph{Leakage and DRAG.}
Because the transmon's $\ket1\to\ket2$ transition is detuned by only
$\alpha=-E_C$, a fast, spectrally broad pulse partially excites
$\ket2$ --- leakage out of the computational subspace, the very effect
\cref{stmt:anharmonicity} warned of. The standard fix, \emph{DRAG}
(derivative removal by adiabatic gate), adds a quadrature component
proportional to $\dot\Omega(t)/\alpha$ that destructively cancels the
off-resonant $\ket2$ amplitude, allowing $\sim20\,$ns gates at the few
hundred MHz anharmonicity of a typical transmon. This is the hardware
realisation of the RWA Rabi problem of \cref{ch:perturbations}, with
the transmon's finite anharmonicity as the leading correction.

\paragraph{Transition.}
Single-qubit rotations plus one entangling gate make a universal set.
The entangling gate is where the qubit--qubit interaction --- and the
transmon's higher levels --- enter.

% --------------------------------------------------------------
\section{Two-qubit gates}
\label{sec:two-qubit-gates}

\begin{phenomenon}
A universal quantum computer needs an entangling two-qubit gate
(\cref{ch:composite}): no amount of single-qubit rotation can entangle.
In circuit QED the entanglement comes from a qubit--qubit interaction
--- either direct capacitive coupling or a virtual exchange through a
shared bus resonator --- which by itself is always on. The art of a
two-qubit gate is to \emph{activate} this interaction for a controlled
duration, either by tuning two qubits into a special resonance with a
flux pulse or by driving them with microwaves, so that an entangling
operation accumulates and then switches off.
\end{phenomenon}

\begin{statement}[\textbf{Three families of two-qubit gate}]
\label{stmt:two-qubit-gates}
Superconducting two-qubit gates fall into three families:
\begin{itemize}
\item \textbf{iSWAP (resonant exchange).} Two qubits coupled by an
  exchange interaction $\hbar J(\hat\sigma_+^{(1)}\hat\sigma_-^{(2)}+\hc)$,
  with $J=g_1g_2/\Delta$ a virtual-photon exchange through a bus, are
  tuned into resonance by a flux pulse. They then swap excitations:
  $\ket{01}\leftrightarrow\ket{10}$ with an $i$ phase. Holding for
  $Jt=\pi/2$ gives the iSWAP; for $Jt=\pi/4$ the entangling
  $\sqrt{\text{iSWAP}}$.
\item \textbf{CZ (controlled phase) via $\ket{11}$--$\ket{02}$.} Tuning
  the two-excitation state $\ket{11}$ near the non-computational
  $\ket{02}$ (a doubly excited transmon, reachable \emph{because} of
  the finite anharmonicity) produces an avoided crossing of size
  $\sqrt2\,g$. Moving adiabatically in and back out makes $\ket{11}$
  acquire a phase the other states do not; when that conditional phase
  reaches $\pi$, the operation is a controlled-$Z$.
\item \textbf{Cross-resonance (all-microwave).} For two fixed-frequency
  qubits with static coupling, driving the \emph{control} qubit at the
  \emph{target} qubit's frequency generates, to leading order, an
  effective $\hat\sigma_z^{(1)}\hat\sigma_x^{(2)}$ ($ZX$) interaction:
  the target rotates in a direction conditioned on the control state,
  which is locally equivalent to a CNOT.
\end{itemize}
Each, combined with single-qubit rotations, is universal.
\end{statement}

\paragraph{From the abstract CNOT to hardware.}
These gates are the concrete realisations of the abstract CNOT that
produced an ebit in \cref{ch:composite}: the controlled-phase and
cross-resonance gates are locally equivalent to CNOT, and entangle an
initially product state into a Bell state exactly as that example
required. The choice among families is architectural: tunable-frequency
processors favour flux-activated iSWAP/CZ (fast, but flux noise adds a
dephasing channel, \cref{stmt:noise-rates}), while fixed-frequency
processors favour the all-microwave cross-resonance gate (quieter, but
slower and more sensitive to frequency crowding). The CZ family is
notable for turning the transmon's ``defect'' --- its low-lying third
level --- into the very mechanism of the gate.

\paragraph{Transition.}
Control delivered, we close the measurement loop: the same dispersive
coupling, now with $g$ and $\chi$ known, performs readout, and a
parametric amplifier built from a second Josephson element makes that
readout quantum-limited.

% --------------------------------------------------------------
\section{Dispersive readout in hardware}
\label{sec:cqed-readout}

\begin{statement}[\textbf{Hardware dispersive readout}]
\label{stmt:cqed-readout}
In the dispersive regime $|\Delta|\gg g$ the transmon--resonator system
realises the dispersive Hamiltonian \cref{eq:dispersive-H} of
\cref{ch:dispersive}, with cavity pull
$\chi=\dfrac{g^2}{\Delta}\dfrac{\alpha}{\Delta+\alpha}$ now built from
the circuit $g$ of \cref{eq:cqed-jc} and the transmon anharmonicity
$\alpha=-E_C$ of \cref{stmt:transmon}. A probe tone at the resonator
frequency acquires a qubit-state-dependent phase
(\cref{stmt:dispersive-readout}); since the resonator is side-coupled
to the feedline, the measured transmission is the notch dip
\cref{eq:notch-s21}, pulled by $\pm\chi$ between the two qubit states.
The signal leaves through the output line as the field $\hat a_{\rm out}$
of the input--output theory (\cref{stmt:input-output}), where it must be
amplified before it can be digitised.
\end{statement}

\paragraph{Why the first amplifier matters.}
The readout signal is a handful of microwave photons. The first
amplifier in the chain dominates the noise of the whole measurement, so
its added noise sets the readout fidelity. A conventional
(transistor) amplifier adds tens of quanta; to reach high-fidelity
\emph{single-shot} readout one needs an amplifier that adds at most the
half-quantum the laws of quantum mechanics permit. That amplifier is
itself a Josephson circuit, which we now build.

% --------------------------------------------------------------
\section{Parametric processes and quantum-limited amplifiers}
\label{sec:parametric}

\begin{phenomenon}
The Josephson cosine is a nonlinear circuit element, and a nonlinearity
driven by a strong \emph{pump} mixes frequencies: energy from the pump
can be converted into pairs of photons in other modes. Expanding the
cosine to cubic order (in a flux-pumped, inversion-asymmetric circuit)
gives three-wave mixing, $\hat a^\dagger\hat a^\dagger\hat a_p$; to
quartic order (the symmetric transmon Kerr term) gives four-wave
mixing. Pumped near twice a mode's frequency, these processes amplify a
weak signal --- and, as the time-reverse of the squeezing interaction
of \cref{stmt:squeezing}, they do so while squeezing the conjugate
quadrature.
\end{phenomenon}

\begin{statement}[\textbf{Parametric amplification and the quantum limit}~\cite{clerk2010}]
\label{stmt:parametric-amp}
A Josephson nonlinearity pumped at frequency $\omega_p$ realises, in
the frame rotating at $\omega_p/2$, the degenerate-parametric
(squeezing) Hamiltonian
\begin{equation}
\label{eq:parametric-hamiltonian}
\op H_{\rm par} = \hbar\xi\,\hat a^\dagger\hat a^\dagger + \hc,
\end{equation}
the squeeze generator of \cref{stmt:squeezing} with pump-set strength
$\xi$. This produces two modes of amplification:
\begin{itemize}
\item \textbf{Phase-sensitive (degenerate):} one signal quadrature is
  amplified with gain $G$ while its conjugate is de-amplified
  (squeezed) by $1/G$. Because no second mode is added, this amplifier
  can be \emph{noiseless} on the amplified quadrature --- it adds no
  quanta.
\item \textbf{Phase-preserving (non-degenerate):} both quadratures (or
  a signal and a separate idler) are amplified equally. Such an
  amplifier \emph{must} add at least half a quantum of noise,
  $n_{\rm add}\ge\tfrac12$, the \emph{standard quantum limit} on linear
  amplification.
\end{itemize}
A Josephson parametric amplifier (JPA) operated near this limit is the
front end of the dispersive-readout chain, making single-shot qubit
measurement possible.
\end{statement}

\paragraph{The same squeezing, run forward.}
\Cref{eq:parametric-hamiltonian} is identical to the squeeze
Hamiltonian that generated squeezed light in \cref{stmt:squeezing}: a
parametric amplifier and a squeezer are the same device, distinguished
only by whether one cares about the amplified quadrature or the
squeezed one. The half-quantum floor on phase-preserving gain is the
amplifier face of the vacuum-noise bound of \cref{stmt:quadratures}, and
the squeezed output of a JPA can in turn be injected into a sensor to
beat the standard quantum limit --- the entanglement- and
squeezing-enhanced metrology developed in \cref{ch:coherent-multiqubit}.

\paragraph{Mixing as a unifying language.}
Three- and four-wave mixing in Josephson elements underlie not only
amplifiers but also frequency converters, parametric couplers (an
alternative, tunable route to the two-qubit gates above), and the
microwave single-photon detectors of \cref{ch:microwave-absorption},
where a four-wave-mixing process converts an incoming signal photon
into a detectable qubit excitation. The Josephson nonlinearity is the
single resource behind all of them.

\paragraph{Transition.}
The amplifier completes the \emph{cavity-based} readout chain. But the
cavity is not mandatory. We close the chapter with the opposite limit:
delete the readout resonator and couple the qubit straight to the
measurement line. The scattering theory that results is the same
input--output bookkeeping of \cref{ch:input-output}, now with a
two-level atom where the linear cavity used to be.

% --------------------------------------------------------------
\section{Resonator-free readout: a qubit coupled directly to the feedline}
\label{sec:waveguide-readout}

\begin{phenomenon}
Dispersive readout interposes a cavity between the qubit and the
measurement line. One can instead delete the cavity and couple the
transmon's dipole \emph{directly} to a transmission line, so that the
qubit is the only circuit element the propagating field sees. A
near-resonant probe is then coherently scattered by the two-level
system: the fraction of the wave that excites the qubit and is
re-emitted interferes with the fraction that passes by, carving a
frequency-dependent \emph{dip} in the transmitted amplitude. A single,
perfectly coupled, coherent emitter can in principle reflect the entire
incident wave. This ``atom in a one-dimensional waveguide'' is the
subject of waveguide QED, and the operating principle of
resonator-free charge-sensitive detectors.
\end{phenomenon}

\begin{statement}[\textbf{Coherent scattering of a qubit on a feedline}~\cite{peropadre2013,wiegand2021}]
\label{stmt:waveguide-readout}
Model the transmission line as a continuum of modes $\hat b_\omega$ and
couple the qubit (transition frequency $\omega_0$) to it linearly,
\begin{equation}
\label{eq:waveguide-hamiltonian}
\op H = \frac{\hbar\omega_0}{2}\hat\sigma_z
 + \hbar\!\int\!\dd\omega\;\omega\,\hat b_\omega^\dagger\hat b_\omega
 + \hbar\!\int\!\dd\omega\;g\,\bigl(\hat b_\omega^\dagger\hat\sigma_- + \hat\sigma_+\hat b_\omega\bigr),
\end{equation}
the two-level version of the input--output Hamiltonian
\cref{eq:io-hamiltonian} with the qubit dipole $\hat\sigma_-$ in place
of the cavity operator $\hat a$. The qubit radiates into the line at
the rate $\Gamma_c=2\pi J g^2$ ($J$ the density of line modes), and the
boundary relation \cref{eq:input-output-relation} becomes
\begin{equation}
\label{eq:waveguide-io}
\hat a_{\rm out}(t)=\hat a_{\rm in}(t) - i\sqrt{\Gamma_c}\,\hat\sigma_-(t).
\end{equation}
Driving on the line with detuning $\Delta=\omega_d-\omega_0$ and Rabi
rate $\Omega$, and allowing non-feedline loss $\Gamma_l$ and pure
dephasing $\Gamma_\phi$, the steady state gives the transmitted and
reflected amplitudes
\begin{align}
\label{eq:waveguide-s21}
S_{21}(\Delta) &= 1 - \frac{\Gamma_c}{2\gamma}\,
   \frac{1 - i\Delta/\gamma}{1+(\Delta/\gamma)^2 + 2\Gamma_n/\gamma},\\
\label{eq:waveguide-s11}
S_{11}(\Delta) &= S_{21}(\Delta)-1,
\end{align}
where $\Gamma_r=\Gamma_c+\Gamma_l$ is the total radiative decay rate,
$\gamma=\tfrac12\Gamma_r+\Gamma_\phi=1/T_2$ the total
coherence-decay rate, and $2\Gamma_n=\Omega^2/\Gamma_r=4P/\hbar\omega_d$
the drive expressed as a photon-interrogation rate ($P$ the on-chip
power). In the weak-probe limit $2\Gamma_n\ll\gamma$ the line saturates
negligibly and \cref{eq:waveguide-s21} is a Lorentzian dip,
\begin{equation}
\label{eq:waveguide-weak}
S_{21}(\Delta)=1-\frac{\Gamma_c}{2\gamma}\,\frac{1}{1+i\Delta/\gamma},
\end{equation}
identical in form to the side-coupled notch resonator
\cref{eq:notch-s21}.
\end{statement}

\paragraph{Derivation.}
In the frame rotating at the drive frequency $\omega_d$ the qubit
obeys $\op H_{\rm eff}=\tfrac{\hbar\Delta}{2}\hat\sigma_z
+\tfrac{\hbar\Omega}{2}\hat\sigma_x$ (the qubit form of the
rotating-wave approximation, \cref{stmt:rwa}), and dissipation enters
through the Lindblad master equation of \cref{stmt:lindblad},
\begin{equation}
\label{eq:waveguide-master}
\dot{\op\rho}=-\tfrac{i}{\hbar}\comm{\op H_{\rm eff}}{\op\rho}
 + \Gamma_r\,\mathcal D[\hat\sigma_-]\op\rho
 + \tfrac{\Gamma_\phi}{2}\,\mathcal D[\hat\sigma_z]\op\rho.
\end{equation}
The resulting optical Bloch equations for the coherence and population
are
\begin{align}
\label{eq:waveguide-bloch}
\langle\dot{\hat\sigma}_\pm\rangle &= (\pm i\Delta-\gamma)\langle\hat\sigma_\pm\rangle
   \mp \tfrac{i\Omega}{2}\langle\hat\sigma_z\rangle,\\
\langle\dot{\hat\sigma}_z\rangle &= -i\Omega\bigl(\langle\hat\sigma_+\rangle-\langle\hat\sigma_-\rangle\bigr)
   -\Gamma_r\bigl(\langle\hat\sigma_z\rangle+1\bigr).
\end{align}
Their decaying transients reproduce the coherence times of
\cref{ch:relaxation-dephasing} --- $T_2=1/\gamma$, $T_1=1/\Gamma_r$,
and $\Gamma_\phi=1/T_2-1/2T_1$ --- so the same $T_1,T_2$ measured on a
conventional qubit are read off here from the scattering linewidth.
Setting $\langle\dot{\hat\sigma}_\pm\rangle=\langle\dot{\hat\sigma}_z\rangle=0$
yields the steady-state coherence
\begin{equation}
\label{eq:waveguide-sigma-minus}
\langle\hat\sigma_-\rangle=-\frac{\Gamma_r}{2}\,
  \frac{\Omega(\Delta+i\gamma)}{\Omega^2\gamma+\Gamma_r(\Delta^2+\gamma^2)},
\end{equation}
which, inserted into the boundary relation \cref{eq:waveguide-io} (with
$\Omega=2g\langle\hat a_{\rm in}\rangle$ relating the drive to the
incident amplitude), gives \cref{eq:waveguide-s21}.

\paragraph{The qubit \emph{is} a notch resonator.}
Comparing the weak-probe transmission \cref{eq:waveguide-weak} with the
side-coupled cavity \cref{eq:notch-s21} term by term identifies the
quality factors
\begin{equation}
\label{eq:waveguide-Q}
Q_c=\frac{\omega_0}{\Gamma_c},\qquad
Q_l=\frac{\omega_0}{2\gamma}=\frac{\omega_0}{2}\,T_2,\qquad
Q_i=\frac{\omega_0}{\Gamma_l+2\Gamma_\phi}.
\end{equation}
The external (coupling) $Q$ is fixed by the radiative rate into the
feedline, the loaded $Q$ by the total coherence time, and the internal
$Q$ by every loss that is not the feedline --- exactly the bookkeeping
of \cref{eq:notch-s21}. A qubit on a transmission line therefore
scatters \emph{precisely} as a linear notch resonator does, with one
decisive difference: a two-level system saturates. Each absorbed photon
must be re-emitted before the next can be scattered, so once the
interrogation rate reaches $2\Gamma_n\sim\gamma$ --- of order one
photon per coherence time --- the dip washes out
(\cref{eq:waveguide-s21}) and the steady-state excitation
\begin{equation}
\label{eq:waveguide-p1}
P_1=\frac{1+\langle\hat\sigma_z\rangle}{2}
   =\frac12\,\frac{2\Gamma_n/\gamma}{1+(\Delta/\gamma)^2+2\Gamma_n/\gamma}
\end{equation}
approaches $\tfrac12$. The linear-resonator description holds only at
the single-photon level; this is why direct-coupled readout is
intrinsically a low-power measurement, with no Josephson amplifier in
between to be saturated.

\paragraph{Readout lives in the equatorial plane.}
Because the measured field is the radiated dipole
\cref{eq:waveguide-io}, the signal is proportional to
$\langle\hat\sigma_-\rangle=\tfrac12(\langle\hat\sigma_x\rangle-i\langle\hat\sigma_y\rangle)$,
the \emph{transverse} coherence --- not the population
$\langle\hat\sigma_z\rangle$ that dispersive readout
(\cref{stmt:dispersive-readout}) projects. A qubit in an energy
eigenstate ($\ket0$ or $\ket1$) radiates no coherent field; only states
with equatorial weight scatter. Direct readout is thus a projection
onto the $XY$ plane of the Bloch sphere rather than a quantum
non-demolition measurement of energy, and the calibration sequences of
\cref{ch:experimental-basics} must be adapted accordingly: a Rabi
measurement reads the emitted amplitude (with
$\Omega^2/\Gamma_r=2\Gamma_n=4P/\hbar\omega_d$ fixing the drive
strength), a Ramsey sequence reads the free-precession envelope
$A(\tau)\propto e^{-\tau/T_2^*}\cos(\Delta\tau+\phi_0)$ directly off the
emitted field, and a $T_1$ sequence needs a final $\pi/2$ pulse to
rotate the surviving population into the plane before the line can
report it.

\begin{application}[\textbf{The charge-sensitive transmon detector (SQUAT)}]
\label{app:squat}
Hold $E_J/E_C$ low enough (of order $E_J/E_C\sim25$) that the
transmon's charge dispersion \cref{eq:charge-dispersion} is not fully
suppressed. The even- and odd-charge-parity bands then remain split by
a measurable $2\chi=2\chi_0\cos(\pi n_g)$, so the scattering resonance
sits at $f_r=f_0\pm\chi$ according to the island parity. A single
continuous feedline tone, fit to \cref{eq:waveguide-s21} to extract
$\Gamma_c$, $\gamma$, and the on-chip drive $\Gamma_n$, then monitors
the parity in real time: every quasiparticle that tunnels across the
junction (\cref{sec:parity-switching}) flips the parity and steps the
transmitted phase or amplitude between two values. Deleting the readout
resonator lets such pixels be packed densely and read continuously at
high bandwidth, the design goal of the superconducting
quasiparticle-amplifying transmon~\cite{magoon2026,fink2024}. Optimal
single-tone parity discrimination matches the dispersion to the
linewidth, $2\chi\sim\gamma$, the direct-coupling analogue of the
$2\chi\sim\kappa$ matching behind dispersive readout
(\cref{stmt:reflection}).
\end{application}

% --------------------------------------------------------------
This chapter gave the abstract cavity-QED toolkit its hardware. A
coupling capacitor produces the Jaynes--Cummings interaction with a
computable $g\propto C_g\sqrt{\omega_r\omega_q}/\sqrt{C_rC_\Sigma}$;
microwave drives implement arbitrary single-qubit rotations (with DRAG
taming the transmon's leakage); and entangling gates come in three
families --- resonant iSWAP, $\ket{11}$--$\ket{02}$ controlled-phase,
and all-microwave cross-resonance --- realising the abstract CNOT of
\cref{ch:composite}. The dispersive coupling performs readout with
$\chi$ now built from circuit parameters, and a Josephson parametric
amplifier, the squeezer of \cref{ch:states-of-light} run forward,
supplies the quantum-limited gain that makes single-shot measurement
possible. Removing the cavity altogether and coupling the qubit
straight to the feedline gives the resonator-free limit, where the
qubit scatters the probe exactly as a notch resonator would --- the
basis of the charge-sensitive detectors of Part~VI.

What remains is to put the apparatus to work: to take a freshly
fabricated chip and find its parameters --- the resonator and qubit
frequencies, the coupling, the $\pi$-pulse amplitude, the coherence
times. That bring-up procedure, the daily reality of operating a
superconducting qubit, is the subject of the final chapter of
Part~IV.

\clearpage
\begin{chaptersummary}

\paragraph{Coupling.}
A coupling capacitor gives
$\op H_{\rm int}=2\beta eV_{\rm zpf}\op n(\hat a+\hat a^\dagger)$,
$\beta=C_g/C_\Sigma$, which in the RWA is Jaynes--Cummings with
$g\propto C_g\sqrt{\omega_r\omega_q}/\sqrt{C_rC_\Sigma}$ --- $g$ from
circuit parameters. Vacuum Rabi, cooperativity, and $\chi=g^2/\Delta$
follow.

\paragraph{Single-qubit gates.}
$\op H_d=\hbar\Omega(t)\cos(\omega_dt+\phi)(\hat b+\hat b^\dagger)$;
on resonance,
$\op H_d^{\rm rot}=\tfrac{\hbar\Omega}{2}(\cos\phi\,\hat\sigma_x+\sin\phi\,\hat\sigma_y)$.
Amplitude sets the angle, phase sets the axis, $\op R_z$ is virtual;
DRAG suppresses $\ket2$ leakage.

\paragraph{Two-qubit gates.}
iSWAP (resonant exchange $J=g_1g_2/\Delta$), CZ (via the
$\ket{11}$--$\ket{02}$ avoided crossing, using the third level), and
cross-resonance ($ZX$, all-microwave, fixed frequency). Each + 1-qubit
rotations is universal; all realise CNOT.

\paragraph{Readout and amplification.}
Dispersive readout with $\chi=(g^2/\Delta)\,\alpha/(\Delta+\alpha)$
from circuit parameters. A Josephson parametric amplifier
($\op H_{\rm par}=\hbar\xi\hat a^{\dagger2}+\hc$, the squeezer run
forward) gives phase-sensitive (noiseless) or phase-preserving
($n_{\rm add}\ge\tfrac12$) gain --- the quantum-limited front end of
single-shot readout.

\paragraph{Resonator-free readout.}
Coupling the qubit directly to the feedline gives coherent scattering
$S_{21}=1-\tfrac{\Gamma_c}{2\gamma}\,(1-i\Delta/\gamma)/[1+(\Delta/\gamma)^2+2\Gamma_n/\gamma]$,
$S_{11}=S_{21}-1$. The qubit acts as a notch resonator with
$Q_c=\omega_0/\Gamma_c$, $Q_l=\omega_0/2\gamma=\tfrac{\omega_0}{2}T_2$,
$Q_i=\omega_0/(\Gamma_l+2\Gamma_\phi)$, but it saturates once
$2\Gamma_n\gtrsim\gamma$. The signal $\propto\langle\hat\sigma_-\rangle$
reads the equatorial plane, not $\langle\hat\sigma_z\rangle$; matched
parity readout takes $2\chi\sim\gamma$ (the SQUAT detector).

\end{chaptersummary}

% --------------------------------------------------------------
\section*{Exercises}
\addcontentsline{toc}{section}{Exercises}

\begin{exercise}[Scaling of the coupling]
From $g\propto C_g\sqrt{\omega_r\omega_q}/\sqrt{C_rC_\Sigma}$, by what
factor does $g$ change if the coupling capacitor $C_g$ is doubled, and
separately if both frequencies are halved? Why is $g$ set at design
time?
\begin{solution}
$g\propto C_g$, so doubling $C_g$ doubles $g$. Halving both
frequencies multiplies $\sqrt{\omega_r\omega_q}$ by $\sqrt{(1/2)(1/2)}=1/2$,
halving $g$. Since $C_g,C_r,C_\Sigma$ are fixed by the chip geometry
and $\omega_r,\omega_q$ by the elements, $g$ is frozen at fabrication
(though $\omega_q$, hence $\Delta$ and $\chi$, can still be tuned with a
SQUID). One designs $g$ for the desired dispersive shift and gate
speeds before the chip is made.
\end{solution}
\end{exercise}

\begin{exercise}[Drive phase sets the rotation axis]
Using \cref{eq:rotating-frame-drive}, write the unitary for a
$\pi/2$-pulse with $\phi=0$ and with $\phi=\pi/2$, and identify the
gates.
\begin{solution}
A pulse of area $\theta=\int\Omega\dd t=\pi/2$ gives
$\op U=\exp[-i(\pi/4)(\cos\phi\,\hat\sigma_x+\sin\phi\,\hat\sigma_y)]$.
For $\phi=0$: $\op U=\exp[-i(\pi/4)\hat\sigma_x]=\op R_x(\pi/2)$, a
$90^\circ$ rotation about $x$. For $\phi=\pi/2$:
$\op U=\op R_y(\pi/2)$. The two differ only by the drive phase --- the
same pulse envelope, shifted by $90^\circ$, rotates about a
perpendicular axis.
\end{solution}
\end{exercise}

\begin{exercise}[Why $\op R_z$ can be virtual]
Explain why a $Z$ rotation need not be a physical pulse, and what it
costs.
\begin{solution}
A $Z$ rotation commutes with the qubit's free evolution and only
changes the \emph{phase reference} of subsequent drives. Advancing the
phase $\phi$ of all later pulses by $\lambda$ is equivalent to having
applied $\op R_z(\lambda)$ before them. So $\op R_z$ is implemented in
software by bookkeeping the drive phase --- it is instantaneous,
exact, and costs no gate time or fidelity. This is why single-qubit
control needs only two physical (equatorial) rotations plus virtual
$Z$.
\end{solution}
\end{exercise}

\begin{exercise}[Leakage and the need for DRAG]
A square $\pi$-pulse of duration $\tau=5\,$ns drives a transmon with
$\alpha/h=-250\,$MHz. Estimate the spectral overlap with the
$\ket1\to\ket2$ transition and argue why DRAG is needed.
\begin{solution}
A $\tau=5\,$ns pulse has spectral width $\sim1/\tau=200\,$MHz, which is
comparable to $|\alpha|/h=250\,$MHz. The drive therefore has
substantial spectral weight at the $\ket1\to\ket2$ frequency (offset by
$|\alpha|$), so it excites $\ket2$ --- leakage of order
$(1/\tau)/|\alpha|\sim200/250\sim O(1)$, unacceptable. DRAG adds a
derivative quadrature $\propto\dot\Omega/\alpha$ engineered to cancel
this off-resonant amplitude, restoring a clean two-level rotation at
$5\,$ns. Without it one must slow the pulse to $\tau\gg1/|\alpha|$.
\end{solution}
\end{exercise}

\begin{exercise}[iSWAP timing]
Two qubits are tuned into resonance with exchange coupling
$J/2\pi=10\,$MHz. How long is a full iSWAP, and a $\sqrt{\text{iSWAP}}$?
\begin{solution}
The exchange swaps $\ket{01}\leftrightarrow\ket{10}$ at rate $J$; a full
iSWAP is $Jt=\pi/2$, i.e.\ $t=\pi/(2J)=1/(4\cdot10\,\text{MHz})=25\,$ns
(using $J=2\pi\cdot10\,$MHz, $t=\pi/(2\cdot2\pi\cdot10^7)=25\,$ns). The
entangling $\sqrt{\text{iSWAP}}$ is half the duration, $t=\pi/(4J)\approx12.5\,$ns.
\end{solution}
\end{exercise}

\begin{exercise}[Why CZ uses the third level]
Explain why a controlled-phase gate exploits the $\ket{02}$ state, and
why a perfect two-level system could not host this particular gate
mechanism.
\begin{solution}
The CZ gate needs a state-dependent phase: only $\ket{11}$ should
acquire an extra phase. Tuning $\ket{11}$ near the avoided crossing with
$\ket{02}$ repels it in energy, so $\ket{11}$ accumulates phase at a
different rate than $\ket{00},\ket{01},\ket{10}$, none of which couple
to a doubly excited state. The $\ket{02}$ level exists only because the
transmon is a multilevel anharmonic oscillator; a strict two-level
system has no second excitation to provide the avoided crossing, so
this mechanism is unavailable. The transmon's weak anharmonicity is
thus essential to the gate, not merely tolerated.
\end{solution}
\end{exercise}

\begin{exercise}[Cross-resonance interaction]
The cross-resonance drive yields an effective
$\op H_{\rm CR}=\hbar\,\Omega_{ZX}\,\hat\sigma_z^{(1)}\hat\sigma_x^{(2)}/2$.
Show that evolving for $\Omega_{ZX}t=\pi/2$ is locally equivalent to a
CNOT.
\begin{solution}
$\exp[-i(\pi/4)\hat\sigma_z^{(1)}\hat\sigma_x^{(2)}]$ rotates qubit~2
about $x$ by $+\pi/2$ when qubit~1 is in $\ket0$ ($\sigma_z=+1$) and by
$-\pi/2$ when qubit~1 is in $\ket1$ --- a control-conditioned $\pi/2$
rotation. Combined with single-qubit $\op R_x(\mp\pi/2)$ on the target
and an $\op R_z$ on the control, this composes to a controlled-$X$
(CNOT): the target flips iff the control is $\ket1$. The two-qubit
content is entirely in the $ZX$ evolution; the rest is local.
\end{solution}
\end{exercise}

\begin{exercise}[The half-photon limit]
Explain physically why a phase-\emph{preserving} linear amplifier must
add $\ge\tfrac12$ quantum of noise while a phase-\emph{sensitive} one
need not.
\begin{solution}
A phase-preserving amplifier reproduces \emph{both} quadratures of the
input with gain. But the two quadratures are conjugate
($\comm{\hat X}{\hat P}\neq0$, \cref{stmt:quadratures}); amplifying both
without adding noise would let one measure both below the vacuum bound,
violating the uncertainty principle. The minimum noise that restores
consistency is half a quantum referred to the input
($n_{\rm add}\ge\tfrac12$). A phase-sensitive amplifier reproduces only
\emph{one} quadrature (squeezing the other), measuring nothing about
the conjugate, so no uncertainty bound is threatened and it can in
principle be noiseless on the amplified quadrature.
\end{solution}
\end{exercise}

\begin{exercise}[Amplifier and squeezer are one device]
Identify the parametric Hamiltonian \cref{eq:parametric-hamiltonian}
with the squeeze generator of \cref{stmt:squeezing}, and state what
physical quantity distinguishes ``amplifier'' from ``squeezer''.
\begin{solution}
The squeeze operator is
$\op S(\zeta)=\exp[\tfrac12(\zeta^*\hat a\hat a-\zeta\hat a^\dagger\hat a^\dagger)]$,
generated by exactly $\hat a^{\dagger2}$ and $\hat a^2$ terms --- the
same as $\op H_{\rm par}=\hbar\xi\hat a^{\dagger2}+\hc$. The unitary
evolution it produces stretches one quadrature and squeezes the
conjugate. Called an ``amplifier'' when one reads out the stretched
(amplified) quadrature of an injected signal; called a ``squeezer''
when one uses the squeezed quadrature of the output vacuum. Same
device, same Hamiltonian; the distinction is which quadrature one cares
about.
\end{solution}
\end{exercise}

\begin{exercise}[The Josephson nonlinearity as the common resource]
List the circuit-QED functions in this chapter that ultimately rely on
the Josephson nonlinearity, and explain in one sentence each how.
\begin{solution}
(i) \emph{Qubit:} the cosine's anharmonicity isolates two levels
(\cref{ch:transmon}). (ii) \emph{CZ gate:} the third level it provides
supplies the $\ket{11}$--$\ket{02}$ avoided crossing. (iii)
\emph{Parametric amplifier:} its three-/four-wave mixing converts pump
power into signal gain. (iv) \emph{Parametric couplers and frequency
converters:} the same mixing activates tunable interactions. (v)
\emph{Microwave photon detectors} (\cref{ch:microwave-absorption}):
four-wave mixing converts a signal photon into a qubit excitation. A
single nonlinear, dissipationless element underlies the qubit, its
gates, and its readout amplifier alike.
\end{solution}
\end{exercise}

\begin{exercise}[On-resonance extinction of a qubit on a feedline]
For a qubit coupled directly to a feedline, evaluate the weak-probe
transmission \cref{eq:waveguide-weak} on resonance ($\Delta=0$). What
does $S_{21}(0)$ become for an ideal qubit with no internal loss and no
pure dephasing, and what is the physical interpretation? Express the
on-resonance dip in terms of the quality factors \cref{eq:waveguide-Q}.
\begin{solution}
On resonance \cref{eq:waveguide-weak} gives
$S_{21}(0)=1-\Gamma_c/2\gamma$. For an ideal qubit
$\Gamma_l=\Gamma_\phi=0$, so $\Gamma_r=\Gamma_c$ and
$\gamma=\tfrac12\Gamma_r=\tfrac12\Gamma_c$, hence $S_{21}(0)=1-1=0$:
the qubit reflects the entire incident wave (perfect extinction), the
single-emitter analogue of a lossless, perfectly overcoupled notch
resonator. In general $\Gamma_c/2\gamma=Q_l/Q_c$, so
$S_{21}(0)=1-Q_l/Q_c=Q_i/(Q_i+Q_c)$ --- identical to the notch result
$S_{21}(\omega_c)=Q_i/(Q_i+Q_c)$ of \cref{eq:notch-s21}. Any internal
loss or dephasing ($Q_i<\infty$) lifts the dip off zero, and finite
probe power ($2\Gamma_n>0$) fills it in further through saturation.
\end{solution}
\end{exercise}
